% =============================================================
% Subsection: Brazilian facts (inside Section 2)
% =============================================================

Formal hours are concentrated near the statutory maximum. The baseline maps the DIEESE contracted-hours distribution into three model bins: $8.5\%$ of formal workers in the $\leq 39$h category are assigned to the 36h bin, $26.9\%$ at 40h remain in the 40h bin, and $64.6\%$ in the $\geq 41$h category are assigned to the 44h bin \citep{DIEESE2024}. PNAD Cont\'{i}nua Q4 2024 habitual hours give a flatter distribution, with an average of $41.2$ hours \citep{PNAD2025}. Most formal workers operate at or above 40 hours, the margin a 36-hour cap would bind.

Informality is large and dispersed. The national rate is $37.8\%$ under the narrow PNAD definition and $44.2\%$ under the broad IBGE one \citep{PNAD2025}. The sectoral range runs from $41.2\%$ in services to $71.8\%$ in agriculture (Table~\ref{tab:facts_informality}).

The formal-informal wage premium is positive and moderate. \citet{MeghirNaritaRobin2015} document a conditional wage premium $R\in[1.15, 1.55]$, the ratio of mean formal hourly wages to mean informal hourly wages among otherwise comparable workers. The calibration below uses this premium to pin down the CES substitution elasticity $\sigma_{\text{sub}}$.

Brazilian TFP has been stagnant or declining throughout the post-1990 period. The Penn World Table 11.0 \texttt{rtfpna} series shows an annualized decline of $-0.83\%$ per year from 1990 to 2023 and $-0.82\%$ from 1990 to 2019, before COVID \citep{FeenstraInklaarTimmer2015,PWT2025,FREDPWTBrazil2026}. The best rolling post-1990 decade in the updated series is 2000 to 2010, with only $+0.05\%$ per year. The TFP benchmark in Section~\ref{sec:results} uses this updated series as the anchor for $A_{\text{req}}$.

These facts call for a model with formal workers concentrated near the cap, substitution between formal and informal effective labor, and per-hour productivity that varies with the workweek. The model adds two firm-size groups for incidence and GHH preferences for welfare.

\begin{table}[htbp]
\centering
\caption{Sectoral and national labor-market moments, PNAD Cont\'{i}nua Q4 2024.}\label{tab:facts_informality}
\small
\begin{tabular}{lcccccc}
\hline\hline
Sector & $\lambda_s$ & Informality & Avg.\ hours & $\theta_{36}$ & $\theta_{40}$ & $\theta_{44}$ \\
 & & (broad) & (formal) & & & \\
\hline
Agriculture & 0.076 & 71.8\% & 44.1 & 0.101 & 0.269 & 0.630 \\
Industry    & 0.205 & 44.4\% & 42.5 & 0.048 & 0.417 & 0.535 \\
Services    & 0.720 & 41.2\% & 40.8 & 0.171 & 0.410 & 0.419 \\
\hline
National    & 1.000 & 44.2\% & 41.2 & 0.143 & 0.406 & 0.451 \\
\hline\hline
\end{tabular}
\\[2pt]
\begin{minipage}{0.95\textwidth}
\justifying
{\footnotesize \textit{Notes:} $\lambda_s$ is the sector's share of national employment. Informality uses the broad IBGE definition (VD4009 combined with V4017 CNPJ indicator). ``Avg.\ hours'' is the mean weekly hours of \emph{formal} workers (V4039, habitual). $\theta_{36}$, $\theta_{40}$, $\theta_{44}$ are the shares of formal workers in each contracted-hours bin. \textit{Source:} PNAD Cont\'{i}nua Q4 2024, IBGE microdata, survey-weighted, $N=209{,}311$ \citep{PNAD2025}.}
\end{minipage}
\end{table}
